% =====================================================================
%  PARTIE 3 — Matériel et outils de développement
% =====================================================================
\chapter{Matériel et outils de développement}
\label{ch:materiel}

Les techniques de compression ne suffisent pas à elles seules : leur efficacité
dépend du matériel sur lequel le modèle est exécuté et des outils logiciels qui
assurent la conversion et le déploiement. Cette partie présente trois
architectures représentatives de l'évolution des microcontrôleurs pour le
TinyML, puis la chaîne d'outils logiciels associée.

\section{Architectures matérielles}

\subsection{ESP32-S3 (Espressif)}
\begin{itemize}
  \item \textbf{Cœur} : Xtensa LX7 double cœur à 240~MHz ;
  \item \textbf{Mémoire} : 512~Ko de SRAM, jusqu'à 8~Mo de PSRAM externe ;
  \item \textbf{Accélération} : extensions vectorielles (SIMD) pour les calculs 8/16~bits ;
  \item \textbf{Connectivité} : Wi-Fi et Bluetooth LE intégrés~\cite{espressif-esp32s3}.
\end{itemize}

\subsection{ESP32-P4 (Espressif)}
\begin{itemize}
  \item \textbf{Architecture} : RISC-V double cœur jusqu'à 400~MHz ;
  \item \textbf{Mémoire} : jusqu'à 32~Mo de PSRAM ;
  \item \textbf{Accélération} : extensions dédiées à l'IA, traitement d'image matériel ;
  \item \textbf{Sortie} : 2024~\cite{espressif-esp32p4}.
\end{itemize}

\subsection{STM32N6 (STMicroelectronics)}
\begin{itemize}
  \item \textbf{Cœur} : Arm Cortex-M55 à 800~MHz ;
  \item \textbf{NPU} : accélérateur neuronal Neural-ART à 600~GOPS ;
  \item \textbf{Efficacité} : environ 3~TOPS/W ;
  \item \textbf{Mémoire} : 4,2~Mo de SRAM intégrée~\cite{st-stm32n6}.
\end{itemize}

\begin{table}[H]
\centering
\small
\begin{tabularx}{\linewidth}{>{\bfseries}l X X X}
\toprule
Critère & ESP32-S3 & ESP32-P4 & STM32N6 \\
\midrule
Cœur           & Xtensa LX7, 2 cœurs, 240 MHz & RISC-V, 2 cœurs, 400 MHz & Cortex-M55, 800 MHz \\
SRAM / PSRAM   & 512 Ko / jusqu'à 8 Mo & 768 Ko / jusqu'à 32 Mo & 4,2 Mo SRAM \\
Accélération IA & Instructions SIMD & Extensions IA, traitement d'image & NPU Neural-ART (600 GOPS) \\
Connectivité   & Wi-Fi, BLE & Aucune radio intégrée & Aucune radio intégrée \\
Usage typique  & Audio, capteurs, petite vision & Vision, interface homme-machine & Vision temps réel, modèles plus lourds \\
\bottomrule
\end{tabularx}
\caption{Comparaison des architectures matérielles pour le TinyML.}
\label{tab:materiel}
\end{table}

Ces trois architectures illustrent l'évolution des microcontrôleurs pour le
TinyML, avec une complexité et des performances croissantes : des instructions
vectorielles d'un MCU généraliste jusqu'à l'intégration d'un véritable
accélérateur neuronal. Pour un scénario comme la détection d'anomalies
vibratoires, un ESP32 ou un ESP32-S3 est largement suffisant ; le Wi-Fi
intégré permet en outre de transmettre les alertes.

\section{Chaîne d'outils logiciels}

\begin{figure}[H]
\centering
\begin{tikzpicture}[node distance=0.6cm]
  \node[bloco, minimum width=2.8cm] (a) {1. Entraînement\\\footnotesize(TensorFlow / PyTorch)};
  \node[bloc, minimum width=2.8cm, right=of a] (b) {2. Conversion\\\footnotesize(TFLite Converter)};
  \node[blocv, minimum width=2.8cm, right=of b] (c) {3. Exécution\\\footnotesize(TFLM / ESP-DL)};
  \node[blocv, minimum width=2.8cm, right=of c] (d) {4. Firmware\\\footnotesize(ESP-IDF / CubeIDE)};
  \draw[fleche] (a) -- (b);
  \draw[fleche] (b) -- (c);
  \draw[fleche] (c) -- (d);
\end{tikzpicture}
\caption{Chaîne d'outils logiciels pour le développement TinyML.}
\label{fig:outils}
\end{figure}

\subsection{TensorFlow Lite for Microcontrollers (TFLM)}

TFLM est le moteur d'inférence de référence pour le TinyML~\cite{david2021tflm,tflm-doc} :
\begin{itemize}
  \item interpréteur C++ minimaliste, fonctionnant sans système d'exploitation
        ni allocation dynamique de mémoire ;
  \item empreinte de quelques dizaines de kilooctets seulement ;
  \item modèle intégré sous forme de tableau C dans le firmware ;
  \item \emph{tensor arena} : zone de mémoire statique, réservée à la
        compilation, qui contient les activations pendant l'inférence.
\end{itemize}

La conversion d'un fichier \code{.tflite} en tableau C se fait en une ligne de
commande, puis le modèle est chargé par l'interpréteur (listing~\ref{lst:tflm}).

\begin{lstlisting}[language=C++, caption={Initialisation de l'interpréteur TFLM sur microcontrôleur.}, label={lst:tflm}]
// Généré par : xxd -i modele_int8.tflite > modele_data.cc
#include "modele_data.h"
#include "tensorflow/lite/micro/micro_interpreter.h"
#include "tensorflow/lite/micro/micro_mutable_op_resolver.h"

constexpr int kTensorArenaSize = 20 * 1024;      // 20 Ko de SRAM
alignas(16) static uint8_t tensor_arena[kTensorArenaSize];

const tflite::Model* model = tflite::GetModel(modele_data);
static tflite::MicroMutableOpResolver<3> resolver;
resolver.AddFullyConnected();
resolver.AddRelu();
resolver.AddSoftmax();

static tflite::MicroInterpreter interpreter(model, resolver,
                                            tensor_arena, kTensorArenaSize);
interpreter.AllocateTensors();
\end{lstlisting}

\subsection{Edge Impulse}

Edge Impulse est une plateforme en ligne couvrant l'ensemble du pipeline, de la
collecte des données au déploiement~\cite{edgeimpulse}. Elle permet
d'enregistrer des données depuis la carte, de concevoir le prétraitement
(FFT, MFCC) et le modèle, d'évaluer ses performances estimées sur le
microcontrôleur cible, puis de générer automatiquement une bibliothèque C++ ou
Arduino prête à être compilée.

\subsection{Outils spécifiques aux fabricants}
\begin{itemize}
  \item \textbf{ESP-DL} (Espressif) : bibliothèque optimisée pour les ESP32, qui
        exploite leurs instructions vectorielles~\cite{espressif-espdl} ;
  \item \textbf{STM32Cube.AI} (STMicroelectronics) : outil de conversion et
        d'optimisation des modèles pour les STM32~\cite{st-cubeai} ;
  \item \textbf{NXP eIQ} : environnement de développement pour les
        microcontrôleurs et processeurs NXP~\cite{nxp-eiq}.
\end{itemize}
